\section{Crossing identities and a homogeneous curvature inequality}
\label{sec:curvature}

This section explains the two limiting operations in
\cite[Proposition B.3 and Theorem C.1, pp.~24--28]{Anthropic2026}.
The first removes the number of terms in the crossing identity.
The second converts the resulting first-order length comparison into
the squared-distance inequality of Bruhat and Tits.
We give the details needed for the planar reduction, including the
uniform error bound in the final averaging argument.

\subsection{The crossing comparison}

We use two inputs from \cite{Anthropic2026}.
Proposition A.8, pp.~21--22, permits adjoining random points to a
price: if new labels are randomly assigned to old labels while the
old labels are fixed, a price extension exists in which every new
distance is at most the expected old distance.
The construction preserves all old distances.
Proposition B.2, pp.~22--24, gives an integral decomposition of two
crossing wires indexed by two classes of Frobenius symbols.
A symbol consists of finite subsets $I,J\subseteq\N$ of equal size,
with energy $\sum_{i\in I}(2i+1)+\sum_{j\in J}(2j+1)$.
If $P(E)$ counts symbols of energy at most $E$, Lemma B.1 gives
\[
 \log P(E)\leq \kappa\sqrt E,\qquad
 \log P(E)=(\kappa-o(1))\sqrt E,
\]
with natural logarithms and a fixed $\kappa>0$.
For symbol classes of energies $e_1,e_2$, their crossing tensor is
a sum of $\Theta+1$ coefficient tensors, where
$\Theta=2(e_1+e_2)$.
The coefficient with index $N$ has ranks at most
$(N+1)P(N)$ and $(\Theta-N+1)P(\Theta-N)$ across the two
adjacent-pair partitions.

Here is the exact passage from those statements to prices.
Given $p,q\geq0$, choose energies bounded by
$\lfloor p^2 n\rfloor$ and $\lfloor q^2 n\rfloor$.
There are at most $(E+1)^2$ size-and-energy classes up to energy
$E$, so some classes have logarithmic sizes
$(p+o(1))\kappa\sqrt n$ and $(q+o(1))\kappa\sqrt n$.
For a fixed spectral point, one coefficient tensor has logarithmic
value at least that of the crossing tensor minus
$\log(\Theta+1)$.
After division by $\kappa\sqrt n$ this loss tends to zero.
The two rank bounds give bridge exponents $r,r'\geq0$ satisfying
$r^2+r'^2\leq2(p^2+q^2)$ after passage to a subsequence.
This passage is valid also when an index $N$ stays bounded,
because $\log(N+1)=O(\log n)$ uniformly for $N\leq\Theta$.
Consequently, for points $x,v,y,u,o,o',w,w'$ of a price,
\[
 p\,d(x,v)+q\,d(y,u)
\]
is at most each of
\[
 p\,d(x,o)+q\,d(y,o)+p\,d(v,o')+q\,d(u,o')+r\,d(o,o'),
\]
\[
 p\,d(x,w)+q\,d(u,w)+q\,d(y,w')+p\,d(v,w')+r'\,d(w,w').
\]
The parameters $(r,r')$ lie in a fixed compact set.
Homogeneity in distances therefore passes the assertions from
spectral metrics to every scaled limit defining a price.
This is the first removal of scale; merely taking powers of one
fixed decomposition would not make its number of summands
subexponential in its logarithmic wire size.

Adding the two bounds and using
$r a\leq t r^2/2+a^2/(2t)$ gives, for every $t>0$,
\begin{align}
 2\bigl(p\,d(x,v)+q\,d(y,u)\bigr)
 &\leq p\bigl(d(x,o)+d(v,o')+d(x,w)+d(v,w')\bigr)\notag\\
 &\quad+q\bigl(d(y,o)+d(u,o')+d(u,w)+d(y,w')\bigr)\notag\\
 &\quad+t(p^2+q^2)
       +\frac{d(o,o')^2+d(w,w')^2}{2t}.
 \label{eq:crossing-young}
\end{align}
The crossing comparison is homogeneous in the distances.
In its Young-inequality form, the free parameter $t$ is rescaled
by the same factor as the distances.

\subsection{Translation averages and a parallelogram law}

Suppose that for each positive integer $n$ a price $W_n$ is given
on two copies of $\Z^2$, with
$W_n(P,P+g)\leq K\|g\|_1$ for a fixed $K$.
Average bounded functions over the two boxes
$0\leq i<n^2$, $0\leq k<n$ and take a limit along a fixed free
ultrafilter. Denote this positive linear average by $\mathbb E$.
It is translation invariant: for a fixed translation the fraction
of points lost from either box tends to zero.
Put
\[
 \nu(g)=\bigl(\mathbb E W(P,P+g)^2\bigr)^{1/2},
 \qquad N(g)=\lim_{S\to\infty}\frac{\nu(Sg)}S.
\]
The triangle inequality, Minkowski's inequality, and translation
invariance show that $\nu$ is even and subadditive.
The displayed limit is its subadditive infimum; $N$ is an even,
subadditive, integer-homogeneous seminorm.

Apply \eqref{eq:crossing-young} to translated configurations with
$p=W(P+x,P+v)$ and $q=W(P+y,P+u)$.
Cauchy--Schwarz bounds every mixed average by the product of the
corresponding $\nu$ values, while the squared terms average
exactly. Thus the specialized inequality, with each crossing
weight equal to the corresponding $\nu$-distance, holds for $\nu$.
Replacing all displacement vectors by $S$ times themselves and
dividing by $S^2$ proves it for $N$.

We spell out why this forces the parallelogram identity.
For vectors $a,b,p,q$ define
\[
 \partial p=2N(p)-N(p-a)-N(p-b),\qquad
 \partial' q=2N(q)-N(q-a)-N(q+b).
\]
Use $(x,v,y,u)=(p,-p,q,-q)$ and
$(o,o',w,w')=(a,-a,b,-b)$ in the preceding inequality and
divide by four. It becomes
\[
 N(p)\partial p+N(q)\partial'q
 \leq t\bigl(N(p)^2+N(q)^2\bigr)
       +\frac{N(a)^2+N(b)^2}{2t}.
 \tag{*}
\]
Let $c=a+b$, $e=a-b$, $f=N(c)$, $f'=N(e)$, $M=R^2$, and
$A_i=Mc+ie$, $B_i=Me+ic$, for $0\leq i<R$.
With $\delta_i=N(A_i-a)-N(A_i-2a)$, direct cancellation gives
\[
 \partial A_i+\partial(A_i-a)
 =2\bigl(N(A_i)-N(A_i-c)\bigr)+\delta_i-\delta_{i+1}.
\]
Convexity of $j\mapsto N(jc+ie)$ on integer points gives
\[
 N(A_i)-N(A_i-c)
 \geq \frac{N(Mc+ie)-N(ie)}M
 \geq f-\frac{2if'}M.
\]
The $\delta_i$ are uniformly bounded and their differences
telescope. Also $\partial$ and $\partial'$ are bounded by
$N(a)+N(b)$, and $N(A_i)=Mf+O(R)$.
Apply (*) to $(A_i,B_i)$ and $(A_i-a,B_i-a)$ with
$t=1/(2M)$, add, average over $i$, and divide by $2M$.
The same computation with $b$ replaced by $-b$ treats $B_i$.
It follows that
\[
 f^2+f'^2
 \leq \tfrac12(f^2+f'^2)+N(a)^2+N(b)^2+O(R^{-1}).
\]
Letting $R$ tend to infinity proves one parallelogram inequality.
Applying that inequality to $a+b,a-b$ and using homogeneity
proves its reverse. Hence
\begin{equation}
 N(a+b)^2+N(a-b)^2=2N(a)^2+2N(b)^2 .
 \label{eq:average-parallelogram}
\end{equation}

\subsection{The fan and the midpoint inequality}

Let $x,v,m,z$ be points of a price with
$d(x,m),d(v,m)\leq D$ and $d(x,v)=2D$.
The triangle inequality makes both first distances equal to $D$.
Put $L_*=D+d(z,m)$ and
\[
 \Gamma=d(z,m)^2-\tfrac12d(z,x)^2-\tfrac12d(z,v)^2.
\]
Fix $\theta=1/T$, with an integer $T\geq2$.
On each of two sheets, adjoin $\gamma_i$, $0\leq i\leq n^2$,
by choosing $m$ with probability $i/n^2$ and otherwise $x$,
respectively $v$, using one common uniform random variable.
Then adjoin $G(i,j)$, $0\leq j\leq Tn^2$, by choosing $z$
with probability $j/(Tn^2)$ and otherwise $\gamma_i$,
again with one common uniform random variable.
The random-point extension gives, with $L_i=d(z,\gamma_i)$
and $\rho=1-\theta j/n^2$,
\[
 d(\gamma_i,\gamma_{i+1})\leq D/n^2,\qquad
 d(z,G(i,j))=\rho L_i,
\]
\[
 d(G(i,j),G(i,j+1))\leq \theta L_i/n^2,\qquad
 d(G(i,j),G(i+1,j))\leq \rho D/n^2.
 \tag{**}
\]
The equality follows because the upper bounds for the distances
to $z$ and to $\gamma_i$ add to $L_i$.
In particular $|L_{i+1}-L_i|\leq D/n^2$ and $L_i\leq L_*$.

Map $(i,k)$ to $G(i',j')$, where $i'$ is $i$ clamped to
$[0,n^2]$ and $j'$ is $i+k+n$ clamped to $[0,Tn^2]$.
Pull back $n^2$ times the price. These are the prices $W_n$
used above, with a common Lipschitz constant $2D+d(z,m)$.
Let $\Lambda^2=\mathbb E L_i^2\leq L_*^2$.
The four estimates required are
\begin{align*}
 N(1,0)&\geq D,\\
 N(1,-1)^2&\leq D^2(1-\theta+\theta^2),\\
 N(0,1)^2&\leq\theta^2\Lambda^2,\\
 N(1,-1-A)^2&\geq
 A^2\theta^2\Lambda^2+A\theta\Gamma
                  -4A\theta^2DL_* \qquad(A\in\N).
\end{align*}
For the first, the distance from the initial endpoint increases
by at most the length of each step.
Summing its increments in direction $(1,0)$ telescopes.
The final point is the same on the two sheets, so its distances
from $x$ and $v$ add to at least $2D$; the initial endpoint error
is $O(n^{-1})$. Translation invariance gives the same lower
average $SD$ for steps of length $S$, and Cauchy--Schwarz gives
the required bound for every $\nu(S(1,0))/S$.

The second estimate follows from (**) because a step $(1,-1)$
keeps $j$ fixed. On the averaging boxes, $j\geq i$ and
the mean of $i/n^2$ tends to $1/2$.
The third follows from the $j$-increment bound in (**).

For the fourth, first fix $S$ and then let $n\geq AS$ tend to
infinity. Write $L=L_i$ and $L'=L_{i+S}$, extended constantly
past $n^2$. The distance from $z$ gives the lower bound
\[
 W_n(P,P+S(1,-1-A))
 \geq\bigl|n^2\rho(L'-L)+AS\theta L'\bigr|.
\]
Square, drop the first square, and use
$2L'(L'-L)\geq L'^2-L^2$.
The mean of $L'^2$ tends to $\Lambda^2$, and telescoping gives
\[
 \mathbb E\bigl(n^2(L'^2-L^2)\bigr)=S\Gamma.
\]
On the box, $0\leq1-\rho\leq2\theta$ for all sufficiently
large $n$, while
$n^2|L'^2-L^2|\leq2SDL_*$.
Replacing $\rho$ by one therefore loses at most
$4AS^2\theta^2DL_*$.
Divide by $S^2$ and take the infimum over $S$.
This proves the last estimate with a constant independent of
$A$, $\theta$, and $S$.

By \eqref{eq:average-parallelogram}, the function
\[
 l\longmapsto N(1,-1-l)^2-l^2\theta^2\Lambda^2
\]
is concave on the integers. Its value at $A$ is at most its
value at zero plus $A$ times the difference of its values at
zero and at $-1$.
The four estimates, followed by $A\to\infty$, give
\[
 \theta\Gamma\leq-\theta D^2+
       \theta^2\bigl(D^2+\Lambda^2+4DL_*\bigr).
\]
Finally let $\theta\downarrow0$. We obtain
\begin{equation}
 d(z,m)^2\leq
 \tfrac12d(z,x)^2+\tfrac12d(z,v)^2-\tfrac14d(x,v)^2 .
 \label{eq:checked-cn}
\end{equation}
This reconstructs Theorem C.1 of the cited preprint from its
random-point and crossing inputs. In a Euclidean space it is
an equality. Thus this curvature consequence, including its
equality cases, does not exclude the square or its centre.
